\chapter{Fundamental Equations}

This chapter mixes equations from different levels of physics. Start with the conservation laws and Newton's laws, then return to Maxwell's equations, relativity, and quantum mechanics after the symbols feel familiar. For each item, read the sentence before the formula first; the formula is the compact version of that sentence.

\textbf{Reality Mapping:} These equations are the grammar of physics---they tell us what can and cannot happen in the universe. Conservation laws say that certain quantities (energy, momentum, angular momentum) are the same before and after any process. Maxwell's equations say that electric and magnetic fields are two aspects of one thing, and that light is an electromagnetic wave. Einstein's equations say that mass and energy are interchangeable, and that gravity is the curvature of spacetime. Quantum equations say that particles behave like waves, that you cannot know everything at once, and that the vacuum is seething with activity. Together, these equations describe the rules that govern everything from falling apples to exploding stars.

\begin{enumerate}

\item \textbf{Ampere--Maxwell Law:} Electric currents and changing electric fields both create magnetic fields. This equation also predicts that light is an electromagnetic wave.
\[
\nabla\times\mathbf B=\mu_0\mathbf J+\mu_0\epsilon_0\frac{\partial\mathbf E}{\partial t}
\]

\item \textbf{Conservation of Angular Momentum:} The total angular momentum of an isolated system does not change. This is why ice skaters spin faster when they pull their arms in.
\[
\mathbf L=\text{const}
\]

\item \textbf{Conservation of Energy:} Energy can change form but the total amount stays the same.
\[
\Delta E_{\text{total}}=0
\]

\item \textbf{Conservation of Mass:} Mass is neither created nor destroyed.
\[
\Delta m=0
\]

\item \textbf{Conservation of Momentum:} The total momentum of an isolated system does not change.
\[
\sum\mathbf p_i=\text{const}
\]

\item \textbf{Einstein Field Equations:} The core of general relativity. They say that mass and energy curve spacetime, and curved spacetime tells mass how to move. They predict black holes, gravitational waves, and the expansion of the universe.
\[
R_{\mu\nu}-\frac12Rg_{\mu\nu}+\Lambda g_{\mu\nu}
=\frac{8\pi G}{c^4}T_{\mu\nu}
\]

\item \textbf{Faraday's Law:} A changing magnetic field creates an electric field. This is how electric generators work.
\[
\nabla\times\mathbf E=-\frac{\partial\mathbf B}{\partial t}
\]

\item \textbf{Gauss's Law:} Electric charges produce electric fields. The total electric flux through a closed surface equals the enclosed charge divided by $\epsilon_0$.
\[
\nabla\cdot\mathbf E=\frac{\rho}{\epsilon_0}
\]

\item \textbf{Gauss's Law for Magnetism:} There are no magnetic monopoles---magnetic field lines always form closed loops.
\[
\nabla\cdot\mathbf B=0
\]

\item \textbf{Law of Thermodynamics (First):} Energy is conserved. The heat added to a system goes into increasing its internal energy or doing work.
\[
dU=\delta Q-\delta W
\]

\item \textbf{Law of Thermodynamics (Second):} Entropy (disorder) of an isolated system never decreases. Heat flows from hot to cold, never the other way.
\[
\Delta S\geq 0
\]

\item \textbf{Law of Thermodynamics (Third):} As temperature approaches absolute zero, entropy approaches a minimum value. You can never actually reach absolute zero.
\[
\lim_{T\to 0}S=S_0
\]

\item \textbf{Law of Thermodynamics (Zeroth):} If two objects are both in thermal equilibrium with a third, they are in equilibrium with each other. This defines temperature.
\[
\text{If }A\sim C\text{ and }B\sim C,\text{ then }A\sim B
\]

\item \textbf{Lorentz Transformation:} The equations that replace Newton's Galilean transformations when objects move close to the speed of light. They show that time slows down and lengths contract for moving objects, and they ensure that the speed of light is the same for everyone.
\[
x'=\gamma(x-vt),\qquad
t'=\gamma\left(t-\frac{vx}{c^2}\right)
\]

\item \textbf{Newton's First Law (Inertia):} A body stays at rest or moves in a straight line at constant speed unless a net force acts on it.
\[
\sum\mathbf F=0 \;\Rightarrow\; \mathbf v=\text{const}
\]

\item \textbf{Newton's Law of Universal Gravitation:} Any two masses attract each other with a force that depends on their masses and the distance between them. This single law explains planetary orbits, tides, and the structure of galaxies.
\[
F=\frac{GMm}{r^2}
\]

\item \textbf{Newton's Second Law:} The net force on an object equals its mass times acceleration. This is the central equation of mechanics.
\[
\mathbf F=m\mathbf a
\]

\item \textbf{Newton's Third Law:} Every force has an equal and opposite partner. When you push a wall, the wall pushes back with the same force.
\[
\mathbf F_{12}=-\mathbf F_{21}
\]

\item \textbf{Schr\"odinger Equation:} The central equation of quantum mechanics. It predicts how the wave function $\psi$ of a particle evolves over time. If you know the wave function, you can calculate the probability of finding the particle anywhere.
\[
i\hbar\frac{\partial\psi}{\partial t}=\hat H\psi
\]

\end{enumerate}
